\section{Introduction}\label{sec:intro}

An agent that pays invoices, renews cloud resources or trades for a user needs
authority that is real and bounded. It must be able to read the files it works
on, reach the services it uses and sign transactions, but only those the user
meant. For an agent acting through a Soroban smart account, the controls in
use check that authority at the wrong grain.

\emph{Custody policy engines} guard the signing key. Dfns, for example,
evaluates rules such as transaction amount limits and recipient allow-lists
before it signs~\cite{dfnsPolicies}. Those rules ``only evaluate plain value
transfers'': when the amount or recipient cannot be inferred, as for a contract
call or a raw signature request, they ``fail closed and
trigger''~\cite{dfnsPolicies}, and the engine does not ``introspect raw
serialized signature payloads''~\cite{dfnsSigningPolicies}. Other engines
decode contract calls and write rules over methods and arguments
(\S\ref{sec:related}), but the ones whose documentation we read describe
EVM-style calls, and we found none that evaluates a Soroban authorization
tree, in which one signed call nests the calls it makes for the account.

\emph{Process sandboxes} guard the machine. nono confines an agent process with
Landlock on Linux and Seatbelt on macOS: it restricts file system paths,
process execution and network destinations, down to a domain, a URL path or an
HTTP method and path, and it matches command arguments against exact, prefix or
contains patterns~\cite{nono}. It does not read what a call carries, such as
the amount and recipient inside a transaction that a remote key signs, and the
smart account that finally executes the transaction cannot see the sandbox's
rules at all.

We write an agent's authority as a \emph{grant}: a list of caveats, each a
predicate over actions, where an action is allowed only if every caveat holds.
A grant's holder can do two things with it: act within it, and hand a sub-agent
a copy with more caveats added. Issuing, revoking, rotating a holder's key and
refilling a budget are the operations of a separate \emph{admin document} that
agents never hold. One evaluator decides every check, at two boundaries:
\begin{itemize}
\item In the \textbf{Soroban smart account}, the evaluator runs inside a
  delegated signer. Through CAP-0071 delegated signing~\cite{cap71,cap7101} the
  host hands it every context of the account's authorization, so a call the
  account authorizes cannot hide a disallowed one in its authorization tree.
\item In the \textbf{agent's host}, the same evaluator, compiled to the same
  WebAssembly target, runs as a component behind the host that performs the
  agent's file, network and signing requests. Before its key keeper signs a
  transaction or an authorization entry, the host checks every context the
  signature authorizes against the agent's chain grant (\S\ref{sec:host}).
\end{itemize}
A relay that submits transactions for agents holding no keys runs the same
chain check before a custody key signs and pays for them. Figure~\ref{fig:arch}
shows the three deployments.

\begin{figure*}[t]
\centering
\resizebox{0.8\textwidth}{!}{\input{architecture}}
\caption{One evaluator at every boundary an agent's actions cross. Each
boundary checks the agent's grant with the same evaluator before acting; the
admin document that changes grants is the owner's alone.}
\label{fig:arch}
\Description{An agent sends tool calls to the agent host, relay requests to
the relay, and holder signatures to the smart account; each boundary runs the
grant evaluator before reaching files, the network and the key keeper, the
custody key and channel accounts, or the contracts the account calls.}
\end{figure*}

The adversary is an agent the user already authorized, trying to widen its own
authority (\S\ref{sec:threat}): by editing its grant, by hiding a disallowed
call inside an allowed one, or by working around delegation through deployed
contracts, approvals or a sibling's allowance.

\paragraph{Contributions.}
\begin{enumerate}
\item A grant model with four action vocabularies (chain calls, files,
  network, signing) that share one evaluator, delegation by attenuation only,
  and budgets as allowances charged up the delegation tree, which stay safe
  when children's allowances oversubscribe their parent's
  (\S\ref{sec:grants}).
\item A proof from the deployed Rust back to a readable Lean model
  (\S\ref{sec:proofs}). Charon and Aeneas translate the evaluator to
  Lean~\cite{charon,aeneas,lean4}, and we prove that the translation computes
  the model for every input in every vocabulary and for the budget functions.
  Totality, attenuation, chaining and budget safety therefore hold
  \translationtrust. Attenuation and chaining are immediate in the model; what
  the proof adds is their carry-over to the code and the budget invariant
  under oversubscription. A step bound is proven for a counted model of the
  evaluator, not for the code.
\item Two enforcement boundaries built on one evaluator, with no change to the
  smart account or to OpenZeppelin's account library: a CAP-0071 delegated
  signer for Soroban (\S\ref{sec:soroban}) and a WebAssembly component behind
  an agent host (\S\ref{sec:host}), plus a relay that checks before a custody
  key signs (\S\ref{sec:relay}).
\item A separation of what is proven from what is tested, measured, assumed
  or unproven, with the trusted wiring and its known gaps stated
  (\S\ref{sec:tcb}), and an evaluation (\S\ref{sec:eval}): an attack suite of
  76 escalation attempts in Soroban simulation and at the host, all refused
  with state unchanged, the Soroban attempts again on a local Stellar network
  through the relay, with its throughput and a race of sibling spends, and a
  cross-check of the proven step bound against metered Soroban costs.
\end{enumerate}

\paragraph{Claim status.} A tag after a claim gives its status, recorded in a
claim ledger that the build checks against the repository: \claimtag{proven}
by named Lean theorems, \claimtag{tested} by a named test,
\claimtag{measured} by a recorded result file, \claimtag{assumed} as a stated
assumption, or \claimtag{unproven}. Appendix~\ref{app:claims} lists every
claim with its evidence. Grants control access, not information flow: a grant
decides which actions an agent may take, not where the data it reads goes next
through other allowed actions (\S\ref{sec:limits}).
